\begin{corollary}[Quantitative dependence on intrinsic depth]
\label{cor:intrinsic-depth-quantitative}
Fix $\zeta>0$. There are constants $a_\zeta,b_\zeta>0$, independent of
the intrinsic depth, with the following property. Suppose $f$ has a
width-$C$, depth-$D_0$ tanh realization with unit signed margin and
weights bounded by $R\ge1$, where $D_0\ge2$, and set
$S=2+d+C+R$. For every specified learner depth $L\ge D_0$, the sample
bound
\[
 n\ge\left\lceil S^{a_\zeta D_0}
       \bigl(L+\varepsilon^{-1}+\log(2/\delta)\bigr)^{b_\zeta}
       \right\rceil
\]
suffices for posterior Gibbs error at most $\varepsilon$, with probability
at least $1-\delta$ for every input distribution, simultaneously for all
$N\ge C_L(f)$ and $0<\sigma^2\le N^{-(1+\zeta)}$.
The prior and likelihood are those of
Corollary~\ref{cor:fixed-intrinsic-depth}.

Thus the dependence on intrinsic depth is
$\exp[O_\zeta(D_0\log S)]$, while the dependence on learner depth remains
polynomial with degree independent of $D_0$. In particular, if
$D_0\le c\log S$ for a fixed constant $c$, the target-parameter factor
is at most $\exp[O_{\zeta,c}((\log S)^2)]$.
This is a quasipolynomial bound in $S$, not a single polynomial jointly
in intrinsic depth and target size.
\end{corollary}

\begin{proof}
First suppose $L\ge2D_0$. Use the tapered construction in the preceding proof, with
$H=D_0-1$, $M=CB^H$, and $T=\max\{1,M\|v\|_2^2\}$.
Its explicit definition gives $B\le S^5$, so
$M\le S^{5D_0}$ and $T\le MC R^2\le S^{7D_0}$.
Put $m=L-H$. Since $L\ge2D_0$, one has $L/m\le2$ and
\[
 K_L=2H+(m+15)(2T)^{1/m}=O(LT),\qquad
 Q_*^2\le e^{H+17}(2T)^{L/m}
          \le e^{H+17}(2T)^2=S^{O(D_0)}.
\]

For the remaining range $D_0\le L<2D_0$, do not taper the target.
Instead pad its original realization to
$M_0=\lceil\max\{4,C,2R\}\rceil\le S^2$, and put
$T_0=\max\{1,M_0\|v\|_2^2\}\le S^5$.
Balanced cloning to every $N\ge M_0$ gives
\[
 E_{\rm pre}\le2H S^4,\qquad c^2\le2T_0.
\]
At $L=D_0$, retain the readout. For $L>D_0$, append the same scalar
transmission with $m=L-H\ge2$. These constructions give, respectively,
\[
 K_L\le2H S^4+2T_0,
 \qquad
 K_L\le2H S^4+(m+15)(2T_0)^{1/m}.
\]
In either case $K_L\le19LS^5$. Therefore
\[
 Q_*^2\le(20S^5)^L\le S^{7L}\le S^{14D_0}.
\]
The width-$M_0$ witness has weights bounded by $M_0$: the original
weights satisfy $R\le M_0/2$, and the scalar-channel check from the
preceding proof applies because $T_0\le M_0^4$.
Thus $C_L(f)\le M_0$ in this range.

In both depth ranges, the polynomial in
Theorem~\ref{thm:unrestricted-above-one} has degree and constants
depending only on $\zeta$. Substituting the corresponding bounds gives
the stated sample form above the respective base width $M$ or $M_0$.
For smaller admissible widths, the defining-witness and Gaussian-box
argument in the preceding proof gives a sample bound polynomial in
$d,L,M,\varepsilon^{-1},\log(1/\delta)$ (or with $M_0$ in place of $M$),
with degree independent of $D_0$. Substitution gives the same sample form.
Split the failure probability between the two width ranges and enlarge
$a_\zeta,b_\zeta$ to finish.
\end{proof}
